% GENERATED FILE. Edit canonical lessons, then run npm run generate:books.
% canonical-source-hash: fnv1a64:a9831373883a8ab5
% canonical-lessons: TE-C507-svabhavam, TE-C507-vyaktitvam, TE-C507-mancitanam, TE-C507-vinayam, TE-C507-ahankaram

\chapter{Character}
\label{ch:te-character}
\hlchaptermodality{507}

\begin{chapteropening}
\textbf{By the end of this chapter:} \emph{I can describe a person's nature, goodness, modesty or pride.}

Complete the last lesson of chapter 507: I can describe a person's nature, goodness, modesty or pride.
\end{chapteropening}

\section[\texorpdfstring{\te{స్వభావం} (svabhāvaṁ)}{స్వభావం (svabhāvaṁ)}]{\te{స్వభావం} (svabhāvaṁ) — temperament; nature}

\label{lesson:TE-C507-svabhavam}



\begin{quote}
Before the new one: say the Telugu for peace of mind, then the Telugu for relief.
\end{quote}

\subsection*{You'll want to know: \te{స్వభావం}}
\textbf{\te{స్వభావం}} — \emph{svabhāvaṁ} — \textquotedblleft{}temperament; nature\textquotedblright{}.

\begin{quote}
\te{మా} \te{తాతయ్యది} \te{చాలా} \te{మంచి} \te{స్వభావం}.
\end{quote}

\begin{quote}
\emph{mā tātayyadi cālā mañci svabhāvaṁ.}
\end{quote}

\begin{quote}
My grandfather has a very good nature.
\end{quote}

Whose nature it is: -\te{ది} \te{స్వభావం}, as in \te{మా} \te{తాతయ్యది}.

\begin{cousinweb}
\textbf{\te{స్వభావం}} (\emph{svabhāvaṁ}) is borrowed from Sanskrit \emph{svabhāva}, \textquotedblleft{}one's own state, inborn disposition\textquotedblright{} (Monier-Williams), from \emph{sva}, \textquotedblleft{}own\textquotedblright{}, and \emph{bhāva}, \textquotedblleft{}being\textquotedblright{}.
\end{cousinweb}

\begin{grammarlens}[title={What you've built}]
One more word for describing what you lived and felt.
\end{grammarlens}

\subsection*{Your turn}
\begin{itemize}
\raggedright
  \item \emph{Say it:} \emph{svabhāvaṁ}
  \item \emph{Say it:} \emph{svabhāvaṁ}, once more
  \item \emph{Say it:} \emph{mā tātayyadi cālā mañci svabhāvaṁ.} — My grandfather has a very good nature.
  \item \emph{Recall:} say the Telugu for peace of mind, then the Telugu for relief, then say \emph{svabhāvaṁ} again
\end{itemize}

\begin{tcolorbox}[breakable,colback=teal!4,colframe=teal!35!black,title={Before you move on}]
What does \te{స్వభావం} mean? (\textquotedblleft{}Temperament; nature\textquotedblright{}.) Say it once more.
\end{tcolorbox}
\section[\texorpdfstring{\te{వ్యక్తిత్వం} (vyaktitvaṁ)}{వ్యక్తిత్వం (vyaktitvaṁ)}]{\te{వ్యక్తిత్వం} (vyaktitvaṁ) — personality}

\label{lesson:TE-C507-vyaktitvam}



\begin{quote}
Before the new one: say the Telugu for relief, then the Telugu for temperament.
\end{quote}

\subsection*{You'll want to know: \te{వ్యక్తిత్వం}}
\textbf{\te{వ్యక్తిత్వం}} — \emph{vyaktitvaṁ} — \textquotedblleft{}personality\textquotedblright{}.

\begin{quote}
\te{ఆమె} \te{వ్యక్తిత్వం} \te{అందరికీ} \te{నచ్చుతుంది}.
\end{quote}

\begin{quote}
\emph{āme vyaktitvaṁ andarikī naccutundi.}
\end{quote}

\begin{quote}
Everyone likes her personality.
\end{quote}

Built on \te{వ్యక్తి}, a person, + -\te{త్వం}, -ness.

\begin{cousinweb}
\textbf{\te{వ్యక్తిత్వం}} (\emph{vyaktitvaṁ}) is borrowed from Sanskrit \emph{vyaktitva}, \textquotedblleft{}distinctness, individuality, personality\textquotedblright{} (Monier-Williams), built on \emph{vyakti}, \textquotedblleft{}an individual\textquotedblright{}.
\end{cousinweb}

\begin{grammarlens}[title={What you've built}]
One more word for describing what you lived and felt.
\end{grammarlens}

\subsection*{Your turn}
\begin{itemize}
\raggedright
  \item \emph{Say it:} \emph{vyaktitvaṁ}
  \item \emph{Say it:} \emph{vyaktitvaṁ}, once more
  \item \emph{Say it:} \emph{āme vyaktitvaṁ andarikī naccutundi.} — Everyone likes her personality.
  \item \emph{Recall:} say the Telugu for relief, then the Telugu for temperament, then say \emph{vyaktitvaṁ} again
\end{itemize}

\begin{tcolorbox}[breakable,colback=teal!4,colframe=teal!35!black,title={Before you move on}]
What does \te{వ్యక్తిత్వం} mean? (\textquotedblleft{}Personality\textquotedblright{}.) Say it once more.
\end{tcolorbox}
\section[\texorpdfstring{\te{మంచితనం} (mañcitanaṁ)}{మంచితనం (mañcitanaṁ)}]{\te{మంచితనం} (mañcitanaṁ) — goodness; decency}

\label{lesson:TE-C507-mancitanam}



\begin{quote}
Before the new one: say the Telugu for temperament, then the Telugu for personality.
\end{quote}

\subsection*{You'll want to know: \te{మంచితనం}}
\textbf{\te{మంచితనం}} — \emph{mañcitanaṁ} — \textquotedblleft{}goodness; decency\textquotedblright{}.

\begin{quote}
\te{అతని} \te{మంచితనం} \te{వల్ల} \te{అందరూ} \te{అతన్ని} \te{నమ్ముతారు}.
\end{quote}

\begin{quote}
\emph{atani mañcitanaṁ valla andarū atanni nammutāru.}
\end{quote}

\begin{quote}
Because of his decency, everyone trusts him.
\end{quote}

-\te{తనం} makes an abstract noun: \te{మంచి}, good; \te{మంచితనం}, goodness.

\begin{cousinweb}
\textbf{\te{మంచితనం}} (\emph{mañcitanaṁ}) is built from \te{మంచి}, \textquotedblleft{}good\textquotedblright{}, and -\te{తనం}, the ending that makes abstract nouns; the origin of \te{మంచి} is not settled.
\end{cousinweb}

\begin{grammarlens}[title={What you've built}]
One more word for describing what you lived and felt.
\end{grammarlens}

\subsection*{Your turn}
\begin{itemize}
\raggedright
  \item \emph{Say it:} \emph{mañcitanaṁ}
  \item \emph{Say it:} \emph{mañcitanaṁ}, once more
  \item \emph{Say it:} \emph{atani mañcitanaṁ valla andarū atanni nammutāru.} — Because of his decency, everyone trusts him.
  \item \emph{Recall:} say the Telugu for temperament, then the Telugu for personality, then say \emph{mañcitanaṁ} again
\end{itemize}

\begin{tcolorbox}[breakable,colback=teal!4,colframe=teal!35!black,title={Before you move on}]
What does \te{మంచితనం} mean? (\textquotedblleft{}Goodness; decency\textquotedblright{}.) Say it once more.
\end{tcolorbox}
\section[\texorpdfstring{\te{వినయం} (vinayaṁ)}{వినయం (vinayaṁ)}]{\te{వినయం} (vinayaṁ) — humility; modesty}

\label{lesson:TE-C507-vinayam}



\begin{quote}
Before the new one: say the Telugu for personality, then the Telugu for goodness.
\end{quote}

\subsection*{You'll want to know: \te{వినయం}}
\textbf{\te{వినయం}} — \emph{vinayaṁ} — \textquotedblleft{}humility; modesty\textquotedblright{}.

\begin{quote}
\te{అంత} \te{డబ్బు} \te{ఉన్నా} \te{ఆయనకు} \te{వినయం} \te{ఎక్కువ}.
\end{quote}

\begin{quote}
\emph{anta ḍabbu unnā āyanaku vinayaṁ ekkuva.}
\end{quote}

\begin{quote}
Though he has so much money, he is very humble.
\end{quote}

\te{వినయంగా}, humbly, politely.

\begin{cousinweb}
\textbf{\te{వినయం}} (\emph{vinayaṁ}) is borrowed from Sanskrit \emph{vinaya}, \textquotedblleft{}guidance, good breeding\textquotedblright{} (CDIAL 11767).
\end{cousinweb}

\begin{grammarlens}[title={What you've built}]
One more word for describing what you lived and felt.
\end{grammarlens}

\subsection*{Your turn}
\begin{itemize}
\raggedright
  \item \emph{Say it:} \emph{vinayaṁ}
  \item \emph{Say it:} \emph{vinayaṁ}, once more
  \item \emph{Say it:} \emph{anta ḍabbu unnā āyanaku vinayaṁ ekkuva.} — Though he has so much money, he is very humble.
  \item \emph{Recall:} say the Telugu for personality, then the Telugu for goodness, then say \emph{vinayaṁ} again
\end{itemize}

\begin{tcolorbox}[breakable,colback=teal!4,colframe=teal!35!black,title={Before you move on}]
What does \te{వినయం} mean? (\textquotedblleft{}Humility; modesty\textquotedblright{}.) Say it once more.
\end{tcolorbox}
\section[\texorpdfstring{\te{అహంకారం} (ahaṅkāraṁ)}{అహంకారం (ahaṅkāraṁ)}]{\te{అహంకారం} (ahaṅkāraṁ) — arrogance; ego}

\label{lesson:TE-C507-ahankaram}



\begin{quote}
Before the new one: say the Telugu for goodness, then the Telugu for humility.
\end{quote}

\subsection*{You'll want to know: \te{అహంకారం}}
\textbf{\te{అహంకారం}} — \emph{ahaṅkāraṁ} — \textquotedblleft{}arrogance; ego\textquotedblright{}.

\begin{quote}
\te{డబ్బు} \te{వచ్చాక} \te{అతనికి} \te{అహంకారం} \te{పెరిగింది}.
\end{quote}

\begin{quote}
\emph{ḍabbu vaccāka ataniki ahaṅkāraṁ perigindi.}
\end{quote}

\begin{quote}
Once he had money, his arrogance grew.
\end{quote}

The opposite of \te{వినయం}. \te{అహం}, I, + \te{కారం}, making.

\begin{cousinweb}
\textbf{\te{అహంకారం}} (\emph{ahaṅkāraṁ}) is borrowed from Sanskrit \emph{ahaṁkāra}, \textquotedblleft{}self-consciousness, egotism, pride\textquotedblright{} (CDIAL 991), from \emph{aham}, \textquotedblleft{}I\textquotedblright{}, and \emph{kāra}, \textquotedblleft{}making\textquotedblright{}.
\end{cousinweb}

\begin{grammarlens}[title={What you've built}]
One more word for describing what you lived and felt.
\end{grammarlens}

\subsection*{Your turn}
\begin{itemize}
\raggedright
  \item \emph{Say it:} \emph{ahaṅkāraṁ}
  \item \emph{Say it:} \emph{ahaṅkāraṁ}, once more
  \item \emph{Say it:} \emph{ḍabbu vaccāka ataniki ahaṅkāraṁ perigindi.} — Once he had money, his arrogance grew.
  \item \emph{Recall:} say the Telugu for goodness, then the Telugu for humility, then say \emph{ahaṅkāraṁ} again
\end{itemize}

\begin{tcolorbox}[breakable,colback=teal!4,colframe=teal!35!black,title={Before you move on}]
What does \te{అహంకారం} mean? (\textquotedblleft{}Arrogance; ego\textquotedblright{}.) Say it once more.
\end{tcolorbox}
